\documentclass{article}

\usepackage{../../Style}

\title{Abelian Varieties}
\author{Shing Tak Lam}
\date{Lent 2023\footnote{Based on lectures by Tony Scholl. Last updated \today.}}

\begin{document}

\maketitle

\lec

% No lecture Jan 29

Recall that an \emph{elliptic curve} \(E\) over a field \(k\), which is a proper curve which is also a group. In particular, they are projective (in fact, we can write it as a nonsingular plane cubic), and that the group is commutative. If \(k\) is algebraically closed, then \(E(k)\) is divisible, and for all \(n \in \N\),
\[E[n] = \ker(n\cdot : E(k) \to E(k))\]
is finite. In particular, if \(n\) is invertible in \(k\), then
\[E[n] = \left(\Z/n\Z\right)^2\]
On the other hand, if \(\Char(k) = p > 0\), then either
\[E[p^r] = \Z/p^r\Z\]
for all \(r \ge 1\), or
\[E[p^r] = 0\]
for all \(r \ge 1\).

The \emph{Abel-Jacobi} theorem gives a homomorphism
\[E(k) \to \Cl^0(E) = \frac{\{\text{degree 0 divisors}\}}{\{\text{principal divisors}\}}\]
Given by \(P \mapsto (P) - (0)\). In fact, by Riemann-Roch, this is an isomorphism.

Abelian varieties are ``higher dimensional analogues of an elliptic curve''. That is, they are a proper variety over \(k\) which is a group. If \(X\) is an abelian variety, we will see:
\begin{itemize}
    \item \(X\) is projective and nonsingular,
    \item \(X\) is commutative,
    \item if \(k\) is algebraically closed, then \(X(k)\) is divisible, with
    \[X[n] \cong (\Z/n\Z)^{2g}\]
    where \(g = \dim(X)\), \(n\) invertible in \(k\). If \(\Char(k) = p > 0\), then
    \[X[p^r] = \left(\Z/p^r\Z\right)^t\]
    with \(0 \le t \le g\).
    \item and an analogue of the Abel-Jacobi theorem, there exists another abelian variety \(\widehat X\), with \(\dim(\widehat X) = g\), called the dual abelian variety, and an isomorphism
    \[\widehat X(k) \cong \Pic^0(X) \sse \Pic(X)\]
    where \(\Pic(X)\) is the group of isomorphism classes of line bundles on \(X\), which is in turn, isomorphic to the divisor class group. Moreover, we have a morphism \(X \to \widehat X\), which is surjective with finite kernel. Note that we don't need to have \(X \cong \widehat X\).
\end{itemize}

In contrast to the one-dimensional case, we need a lot of algebraic geometry, especially the cohomology of coherent sheaves to prove these. One reason is that unlike elliptic curves, we don't have nice equations for them.

\textbf{Conventions for this course:} Every ring will be commutative unital, and ring homomorphisms will be unital. If \(f : A \to B\) is a ring homomorphism, we say \(B\) is an \(A\)-algebra. We'll write \(ab\) for \(f(a)b\). An \(A\)-module is \emph{finite} if it is finitely generated. An \(A\)-algebra \(B\) is \emph{finite} if it is finite as an \(A\)-module, and it is of \emph{finite type} if it is finitely generated as an \(A\)-algebra.

A \emph{family} \((x_i)_{i \in I}\) of elements of a set \(S\) is just \(x_i \in S\), indexed by a set \(I\). This is the same as a function \(I \to S\). We'll write \(\#S\) for the size of a set \(S\), \(\subset\) and \(\sse\) will be used interchangably. Finally,
\[\N = \{0, 1, \dots\}\]

\section{K\"ahler differentials}

Let \(\varphi : A \to B\) be a ring homomorphism. We'll define a \(B\)-module \(\Omega_{B/A}\), which is the module of (K\"ahler) differentials, with an additive map \(\dd = \dd_{B/A} : B \to \Omega_{B/A}\), such that
\begin{align*}
    \dd(b_1b_2) &= b_1\dd b_2 + b_2 \dd b_1 \\
    \dd a = \dd(\varphi(a)) &= 0 \text{ for all }a \in A 
\end{align*} 
In particular, \(\dd\) is an \(A\)-linear map.

\begin{definition}
    We define \(\Omega_{B/A} = \mcP/\mcQ\), where
    \begin{align*}
        \mcP &= \text{free \(B\)-module on symbols }[b]\text{ for }b \in B \\
        \mcQ &= \text{submodule generated by }[a], [b_1 + b_2] - [b_1] - [b_2], [b_1b_2] - b_1[b_2] - b_2[b_1]
    \end{align*}
    and
    \[\dd b = [b] \mod \mcQ\]
\end{definition}

An \emph{\(A\)-derivation of \(B\) into a \(B\)-module \(M\)} is an additive map \(\DD : B \to M\), such that
\begin{align*}
    \DD(A) &= 0 \\
    \DD(b_1b_2) &= b_1D(b_2) + b_2 D(b_1)
\end{align*}
We will write \(\Der_A(B, M)\) for the set of \(A\)-derivations \(B \to M\). This is a \(B\)-module, with \((b'\DD)(b) = b'\DD(b)\). Moreover, \(\Der_A(B, \cdot)\) is a functor. That is, if we have a \(B\)-module homomorphism \(f : M \to M'\), then we have a map
\[\Der(A, M) \to \Der(A, M')\]
given by \(\DD \mapsto f \circ \DD\). In particular, \((\Omega_{B/A}, \dd_{B/A})\) is the \emph{universal} derivation.

\begin{proposition}
    Suppose \(M\) is a \(B\)-module, then there is an isomorphism of \(B\)-modules,
    \begin{align*}
        \Hom_B(\Omega_{B/A}, M) &\cong \Der_A(B, M) \\
        \psi &\mapsto \psi \circ \dd_{B/A}
    \end{align*}
\end{proposition}

\begin{proof}
    The fact that it is a \(B\)-module homomorphism is trivial. Let \(\DD \in \Der_A(B, M)\). Set \(\widetilde\psi : \mcP \to M\) to be \(B\)-linear, with \(\widetilde\psi([b]) = D(b)\). Since \(D\) is a derivation, \(\widetilde\psi(Q) = 0\), so we have a map \(\psi : \mcP/\mcQ = \Omega_{B/A} \to M\), with \(\psi \circ \dd = D\).

    On the other hand, if \(\psi \circ \dd = 0\), then \(\psi = 0\) since \(\Omega_{B/A}\) is generated by the image of \(\dd\).
\end{proof}

\lec

Another characterisation of \(\Omega_{B/A}\) is as follows. Consider the map \(\mu : B \otimes_A B \to B\), given by \(\mu(b_1 \otimes b_2) = b_1b_2\). In particular, \(\mu\) is a \(B\)-algebra homomorphism, for each of the \(B\)-algebra structues on \(B \otimes_A B\). Let \(J = \ker(\mu)\). Then \(J/J^2\) is a \(B \otimes_A B\)-module, killed by \(J\), so it is a \(B\)-module. In particular, all of the \(B\)-module structures on it agree.

Let
\begin{align*}
    \dd' : B &\to J/J^2 \\
    b &\mapsto (1 \otimes b - b \otimes 1) \mod J^2
\end{align*}
Then \(\dd'\) is an \(A\)-derivation. To see this,
\[\dd'(a) = 1 \otimes a - a \otimes 1 = 0\]
for all \(a \in A\). Next,
\begin{align*}
    b_1\dd'(b_2) &= (b_1 \otimes 1)(1 \otimes b_2 - b_2 \otimes 1) \mod J^2 \\
    &= b_1 \otimes b_2 -  b_1b_2 \otimes 1 \mod J^2
\end{align*}
and by symmetry,
\[b_2\dd'(b_1) = 1 \otimes b_1b_2 - b_1 \otimes b_2 \mod J^2\]
Adding these together gives the Leibniz rule. By the universal property, there exists a unique \(B\)-linear map \(\psi : \Omega_{B/A} \to J/J^2\), such that \(\psi \circ \dd = \dd'\).

\begin{proposition}
    \(\psi\) is an isomorphism.
\end{proposition}

So we could have defined \((\Omega_{B/A}, \dd)\) as \((J/J^2, \dd')\) instead.

\begin{proof}
    Consider the map
    \begin{align*}
        \varphi : B \otimes_A B &\to \Omega_{B/A} \\
        \varphi(b_1 \otimes b_2) &= b_1 \dd(b_2)
    \end{align*}
    Now
    \[J = \left\{\sum b_i \otimes b_i' \ \big\vert\ \sum b_ib_i' = 0\right\}\]
    We claim that \(J\) is generated (as a \(B\)-module) by elements of the form \(1 \otimes b - b \otimes 1\). To see this,
    \[\sum b_i \otimes b_i' = \sum (b_i \otimes 1)(1 \otimes b_i' - b_i' \otimes 1) + \sum b_ib_i' \otimes 1\]
    and the last term is zero. So it is a \(B\)-module under the \(b \otimes 1\) action. Now
    \begin{align*}
        \varphi((1 \otimes b - b \otimes 1)(1 \otimes b' - b' \otimes 1)) = 0
    \end{align*}
    and so \(\varphi\) vanishes on \(J^2\), and it factors through \((B \otimes_AB)/J^2\). It remains to check that \(\varphi \circ \psi = \id\) and \(\psi \circ \varphi\vert_{J/J^2} = \id\). We leave this as an exercise.
\end{proof}

\begin{remark}
    Let \(C = (B \otimes_A B)/J^2\), then we have
    % https://q.uiver.app/#q=WzAsNCxbMCwwLCJCIl0sWzEsMCwiQiBcXG90aW1lc19BQSJdLFsyLDAsIkMiXSxbMiwxLCJCIl0sWzAsMSwiXFxjb25nIl0sWzEsMiwiIiwwLHsic3R5bGUiOnsidGFpbCI6eyJuYW1lIjoiaG9vayIsInNpZGUiOiJ0b3AifX19XSxbMiwzLCJcXG11Il0sWzAsMywiXFxpZCIsMl1d
    \[\begin{tikzcd}
        B & {B \otimes_AA} & C \\
        && B
        \arrow["\cong", from=1-1, to=1-2]
        \arrow[hook, from=1-2, to=1-3]
        \arrow["\mu", from=1-3, to=2-3]
        \arrow["\id"', from=1-1, to=2-3]
    \end{tikzcd}\]
    This gives us an isomorphism of \(B = B \otimes_A A\)-algebras,
    \[B \oplus J/J^2 \cong C\]
    The ring structure on the left hand side is given by
    \[(b, f)\cdot (b', f') = (bb', bf' + b'f)\]
\end{remark}

From the definition, any commutative square of rings
% https://q.uiver.app/#q=WzAsNCxbMCwwLCJCIl0sWzAsMSwiQSJdLFsxLDEsIkEnIl0sWzEsMCwiQiciXSxbMCwzXSxbMSwwXSxbMSwyXSxbMiwzXV0=
\[\begin{tikzcd}
	B & {B'} \\
	A & {A'}
	\arrow[from=1-1, to=1-2]
	\arrow[from=2-1, to=1-1]
	\arrow[from=2-1, to=2-2]
	\arrow[from=2-2, to=1-2]
\end{tikzcd}\]
induces a map
\[\Omega_{B/A} \to \Omega_{B'/A'}\]
which is \(B\)-linear, commuting with \(\dd\). Moreover, this is transitive for \(B \to B' \to B''\).

\begin{proposition}
    \begin{enumerate}[(i)]
        \item 
        If \(B' = B \otimes_A A'\), then
        \[\Omega_{B/A} \otimes_A A' \cong \Omega_{B/A} \otimes_B B' \cong \Omega_{B'/A'}\]
        \item if \(A = A'\), and \(B' = S^{-1}B\) is a localisation of \(B\), then
        \[\Omega_{S^{-1}B/A} = S^{-1}\Omega_{B/A}\]
    \end{enumerate}
\end{proposition}

\begin{proof}
    Exercise.
\end{proof}

\begin{example}
    If \(B = A[t_1, \dots, t_n]\) is the polynomial algebra, then
    \[\Omega_{B/A} = \oplus_{i=1}^n B\dd t_i\]
    is free. To see this,
    \[B \otimes_A B = A[\{t_i \otimes 1, 1 \otimes t_j\}] = A[\{t_i \otimes 1, z_i\}] = B[\{z_i\}]\]
    where \(z_i = 1 \otimes t_i - t_i \otimes 1\). Using this, the map \(\mu : B \otimes_A B \to B\) is just \(z_i \mapsto 0\) and the identity map on \(B\). So we have that
    \[J = \inner{z_1, \dots, z_n}\]
    and
    \[J^2 = \inner{z_iz_j}\]
    Thus, \(J/J^2\) is the free module
    \[\bigoplus_i B \cdot (z_i \mod J^2)\]
    But \(z_i \mod J^2\) is just \(\dd'(z_i) = \dd t_i\).
\end{example}

We have two exact sequences for \(\Omega\):

\begin{proposition}
    [transitivity] if \(A \to B \to C\) are ring homomorphisms, then
    % https://q.uiver.app/#q=WzAsNCxbMCwwLCJcXE9tZWdhX3tCL0F9XFxvdGltZXNfQkMiXSxbMSwwLCJcXE9tZWdhX3tDL0F9Il0sWzIsMCwiXFxPbWVnYV97Qy9CfSJdLFszLDAsIjAiXSxbMCwxXSxbMSwyXSxbMiwzXV0=
\[\begin{tikzcd}
	{\Omega_{B/A}\otimes_BC} & {\Omega_{C/A}} & {\Omega_{C/B}} & 0
	\arrow[from=1-1, to=1-2]
	\arrow[from=1-2, to=1-3]
	\arrow[from=1-3, to=1-4]
\end{tikzcd}\]
is exact.
\end{proposition}

\begin{proposition}
    [second exact sequence] If \(A \to B \surj C\) are ring homomorphisms, with say \(C = B/I\), then
    % https://q.uiver.app/#q=WzAsNCxbMSwwLCJcXE9tZWdhX3tCL0F9XFxvdGltZXMgQyJdLFsyLDAsIlxcT21lZ2Ffe0MvQX0iXSxbMywwLCIwIl0sWzAsMCwiSS9JXjIiXSxbMCwxXSxbMSwyXSxbMywwXV0=
\[\begin{tikzcd}
	{I/I^2} & {\Omega_{B/A}\otimes C} & {\Omega_{C/A}} & 0
	\arrow[from=1-2, to=1-3]
	\arrow[from=1-3, to=1-4]
	\arrow[from=1-1, to=1-2]
\end{tikzcd}\]
    is exact, where \(f \mod I^2 \mapsto \dd f \otimes 1\). This is an exact sequence of \(C\)-modules.
\end{proposition}

Proof of both have been left as exercises.

\begin{corollary}
    If \(B = A[\{x_i\}]\), and \(C = B/\inner{\{f_j\}}\), then
    \[\Omega_{C/A} = \frac{\bigoplus_i C \dd x_i}{\sum C \dd f_j}\]
    In this case,
    \[\dd f_j = \sum_i \pdv{f_j}{x_i}\dd x_i\]
\end{corollary}

\subsection{Sheafification}

Let \(f : X \to Y\) be a morphism of schemes. We would like to define a quasicoherent \(\mcO_X\)-module \(\Omega_{X/Y}\), and a map
\[\dd : \mcO_X \to \Omega_{X/Y}\]
which is \(f^{-1}\mcO_Y\)-linear, and satisfies the Leibniz rule for sections. This is called the \emph{sheaf of relative differentials}, or the \emph{cotangent sheaf}. There are two constructions, let's see the `dirty' one first.

If we have open affines \(U = \Spec(B) \sse X\), \(V = \Spec(A) \sse Y\), with \(f(U) \sse V\), then we can define
\[\Omega_{X/Y}\vert_U = \widetilde{\Omega_{B/A}}\]
where \(\widetilde{\Omega_{B/A}}\) is the quasicoherent \(\mcO_U\)-module attached to the \(B\)-module \(\Omega_{B/A}\). The map \(\dd\) is induced by
% https://q.uiver.app/#q=WzAsNCxbMCwwLCJcXG1jT19YKFUpIl0sWzEsMCwiXFxPbWVnYV97WC9ZfShVKSJdLFswLDEsIkIiXSxbMSwxLCJcXE9tZWdhX3tCL0F9Il0sWzIsMywiXFxkZCIsMl0sWzAsMiwiXFxjb25nIiwyXSxbMSwzLCJcXGNvbmciXSxbMCwxLCJcXGRkIl1d
\[\begin{tikzcd}
	{\mcO_X(U)} & {\Omega_{X/Y}(U)} \\
	B & {\Omega_{B/A}}
	\arrow["\dd"', from=2-1, to=2-2]
	\arrow["\cong"', from=1-1, to=2-1]
	\arrow["\cong", from=1-2, to=2-2]
	\arrow["\dd", from=1-1, to=1-2]
\end{tikzcd}\]
Functoriality means that we can glue to get a quasicoherent sheaf on \(X\). We will omit the details.

For the second construction, consider the diagonal map \(\Delta : X \to X \times_Y X\). This map is an immersion, and so it factors as
% https://q.uiver.app/#q=WzAsMyxbMCwwLCJYIl0sWzEsMSwiVSJdLFsyLDAsIlhcXHRpbWVzX1lYIl0sWzAsMSwiXFx0ZXh0e2Nsb3NlZH0iLDJdLFsxLDIsIlxcdGV4dHtvcGVufSIsMix7InN0eWxlIjp7InRhaWwiOnsibmFtZSI6Imhvb2siLCJzaWRlIjoidG9wIn19fV0sWzAsMl1d
\[\begin{tikzcd}
	X && {X\times_YX} \\
	& U
	\arrow["{\text{closed}}"', from=1-1, to=2-2]
	\arrow["{\text{open}}"', hook, from=2-2, to=1-3]
	\arrow[from=1-1, to=1-3]
\end{tikzcd}\]
Let \(i : X \inj U\) be the closed immersion. Define
\[\Omega_{X/Y} = i^*\left(\mcI_{X/U}/\mcI_{X/U}^2\right)\]
where \(\mcI_{X/U}\) is the ideal sheaf. When \(X = \Spec(B)\) and \(Y = \Spec(A)\), this is just \(\widetilde{\Omega_{B/A}}\).
\end{document}